%% sections/01_introduction.tex

\section{Introduction}
\label{sec:intro}

Long-term numerical integration of asteroid orbits requires
integrators that are simultaneously accurate, fast, and structurally
stable. Symplectic integrators \citep{WisdomHolman1991, Yoshida1990}
are the method of choice because they preserve the geometric structure
of Hamiltonian flow, preventing secular energy drift over millions of
orbital periods \citep{HairerLubichWanner2006}. However, standard
adaptive integrators with symplectic kernel use a fixed time step, which is computationally
inefficient for eccentric orbits: the step appropriate near perihelion
is unnecessarily small near aphelion.

\AAS\ is part of the \astdyn\ C++20 library, whose primary use
case is the prediction of stellar occultations by solar-system
minor bodies for the \textit{ITALOccult} network of amateur and
semi-professional observers.
The target audience is therefore twofold: (i)~researchers and
advanced practitioners who require a validated, open-source
propagator with analytical STM support for orbit determination
pipelines; and (ii)~community-based observers who need accurate,
reproducible ephemerides on short timescales (hours to days) for
occultation planning. These two use cases impose complementary
requirements: long-term energy conservation for the former,
and short-term positional accuracy ($\lesssim 1$\,mas) for
the latter.

Adaptive time-step methods for adaptive integrators with symplectic kernel were pioneered
by \citet{PretoTremaine1999} and \citet{MikkolaTanikawa1999}, who
showed that a Sundman-type reparametrization $dt = g(\mathbf{q})\,ds$
can preserve symplecticity when $g$ depends only on position
coordinates. \citet{BlanesIserles2012} subsequently unified and
extended these approaches in a general framework for separable
Hamiltonians.
The key design question is the choice of $g(\mathbf{q})$. Previous
implementations have used functions of inter-particle separation
\citep{PretoTremaine1999} or the logarithmic Hamiltonian
\citep{MikkolaTanikawa1999}. High-order fixed-step integrators
such as \saba\ \citep{LaskarRobutel2001} provide excellent long-term
accuracy for nearly circular orbits, but cannot adapt to the
multiscale dynamics of high-eccentricity objects.

We propose a physically motivated step function based on the spectral
radius of the gravitational potential Hessian
$\rhoH = \rho(\nabla^2\Phi(\mathbf{q}))$.
This quantity: (i) is analytically computable for Keplerian plus
$J_2$ perturbations; (ii) directly measures local orbital curvature;
(iii) is already evaluated during force computation, adding no overhead.
For a Keplerian orbit the resulting step $\Delta t \propto r^{3/2}$
is identical in scaling to the optimal step derived by
\citet{PretoTremaine1999} from a time-transformation analysis, here
obtained from a purely geometric argument.

A further innovation is the analytical propagation of the State
Transition Matrix (STM) within the Drift-Kick-Drift sequence via
the same Hessian. Standard numerical STM methods require 12
additional integrations per step \citep[\S4.2]{bigi_astdyn_2025};
the \AAS\ approach reduces this to a single matrix update per sub-step.

The \AAS\ integrator is embedded in the \astdyn\ C++20 library
\citep{bigi_astdyn_2025}, which provides a complete pipeline for
high-precision asteroid orbit determination and stellar occultation
prediction. 
A companion paper (Bigi, in prep.) describes the full
library pipeline from MPC observations to stellar occultation
path predictions, validated against JPL Horizons with
residuals $< 2.5\,\mu$m($10^{-9}$\,AU).

Table~\ref{tab:tools} compares \astdyn\ with existing
orbit propagation frameworks. Unlike MONTE \citep{evans2018},
GODOT \citep{godot2020}, and Tudatpy \citep{dirkx2022},
which target spacecraft mission design, \astdyn\ is optimised
for the specific requirements of ground-based occultation
prediction: sub-millisecond latency per query, analytical STM
for covariance propagation, and direct integration with the
MPC/Horizons data pipeline. Unlike OrbFit \citep{milani2010}
and Find\_Orb, \astdyn\ provides a modern C++20 API with
full uncertainty propagation and is openly licensed.

\begin{table*}
\centering
\caption{Comparison of orbit propagation tools.
\label{tab:tools}}
\begin{tabular}{lcccc}
\toprule
Tool & Language & Open source & Analytical STM & Primary use \\
\midrule
MONTE        & Python/C++ & No  & Yes & Spacecraft nav \\
GODOT        & Python     & No  & Yes & Spacecraft nav \\
Tudatpy      & Python/C++ & Yes & Yes & Spacecraft nav \\
OrbFit       & Fortran    & Yes & No  & Asteroid OD    \\
Find\_Orb    & C++        & Yes & No  & Asteroid OD    \\
REBOUND      & C/Python   & Yes & No  & N-body sim     \\
\textbf{\astdyn} & \textbf{C++20} & \textbf{Yes} &
  \textbf{Yes} & \textbf{Occultations} \\
\bottomrule
\end{tabular}

\medskip
\parbox{0.95\textwidth}{\footnotesize
OD = orbit determination. STM = State Transition Matrix.
\astdyn\ is the only open-source tool combining analytical STM with an
occultation-oriented pipeline.
}
\end{table*}
Section~\ref{sec:discussion} discusses the results and limitations.
Section~\ref{sec:conclusions} summarizes the conclusions.
